% Block 1 of lista-reais-lectures-01-05 (input by the drivers in Listas/).
% Order and ratings: lista-reais-lectures-01-05/order.md.
\bloco{1}{Static games I: dominance and iterated elimination}{Lecture 1}

% ---------------------------------------------------------------------------
\exercicio[1]{Iterated Elimination}{Foundations}{Tadelis (2013), \emph{Game Theory: An
Introduction}, Ch.~4 ``Rationality and Common Knowledge'', Exercise 4.5 ``Iterated
Elimination'', p.~77 [PDF 96]. No answer in the book. Reproduced verbatim.}

\rot{4.5}\textbf{Iterated Elimination:} In the following normal-form game, which strategy profiles
survive iterated elimination of strictly dominated strategies?
\jogo{Player 1}{Player 2}{\begin{tabular}{c|c|c|c|}
\hd{} & \hd{$L$} & \hd{$C$} & \hd{$R$}\\\cline{2-4}
$U$ & 6, 8 & 2, 6 & 8, 2\\\cline{2-4}
$M$ & 8, 2 & 4, 4 & 9, 5\\\cline{2-4}
$D$ & 8, 10 & 4, 6 & 6, 7\\\cline{2-4}
\end{tabular}}
\resolespaco

% ---------------------------------------------------------------------------
\exercicio[1]{Iterated Deletion of Strictly Dominated Strategies (IDSDS)}{Foundations}{Muñoz-Garcia \&
Toro-Gonzalez (2019), \emph{Strategy and Game Theory: Practice Exercises with Answers}, 2nd ed.,
Ch.~1 ``Dominance Solvable Games'', Exercise 13 ``Iterated Deletion of Strictly Dominated Strategies
(IDSDS)''$^{\mathrm{A}}$, p.~26 [PDF 39]. Worked answer in the book, right after the statement.
Reproduced verbatim (statement only); Fig.~1.33, printed inside the answer on the same page,
rebuilt.}

Consider the following normal form game (Fig.~1.33)
\begin{enumerate}[label=(\alph*),leftmargin=2em]
\item Find strictly dominant strategies (if any) for player 1 and for player 2.
\item Find strictly dominated strategies (if any) for player 1 and for player 2.
\item If you apply iterative deletion of strictly dominated strategies (IDSDS), what is the surviving
strategy pair (or pairs)? Explain the steps you use in IDSDS, and why you use them.
\end{enumerate}
\jogo{Player 1}{Player 2}{\begin{tabular}{c|c|c|}
\hd{} & \hd{\emph{Left}} & \hd{\emph{Right}}\\\cline{2-3}
\emph{Up} & -10, 10 & 0, 12\\\cline{2-3}
\emph{Middle} & -12, 0 & 2, 2\\\cline{2-3}
\emph{Down} & -9, 0 & 1, 0\\\cline{2-3}
\end{tabular}}
\vspace{-6pt}\begin{center}\small\textit{Fig.~1.33} Normal-form game\end{center}
\resolespaco

% ---------------------------------------------------------------------------
\exercicio[2]{IESDS and pure-strategy Nash equilibria}{Foundations}{Gibbons (1992), \emph{Game Theory
for Applied Economists}, Ch.~1 ``Static Games of Complete Information'', \S1.5 Problems
(Section 1.1), Problem 1.2, p.~48 [PDF 56]. No answer in the book. Reproduced verbatim; the payoff
pair in cell $(T,L)$ is printed ``2.0'' (comma typo) and is typeset ``2, 0''.}

\rot{1.2.}In the following normal-form game, what strategies survive iterated elimination of
strictly dominated strategies? What are the pure-strategy Nash equilibria?
\begin{center}\begin{tabular}{c|c|c|c|}
\hd{} & \hd{$L$} & \hd{$C$} & \hd{$R$}\\\cline{2-4}
$T$ & 2, 0 & 1, 1 & 4, 2\\\cline{2-4}
$M$ & 3, 4 & 1, 2 & 2, 3\\\cline{2-4}
$B$ & 1, 3 & 0, 2 & 3, 0\\\cline{2-4}
\end{tabular}\end{center}
\resolespaco

% ---------------------------------------------------------------------------
\exercicio[2]{Prove Proposition 4.1}{Exam level}{Tadelis (2013), \emph{Game Theory: An Introduction},
Ch.~4 ``Rationality and Common Knowledge'', Exercise 4.1, p.~76 [PDF 95]. No answer in the book.
Reproduced verbatim. The exercise states the proposition it asks to prove.}

\rot{4.1}\textbf{Prove Proposition 4.1:} If the game
$\Gamma = \langle N, \{S_i\}_{i=1}^{n}, \{v_i\}_{i=1}^{n} \rangle$ has a strictly dominant strategy
equilibrium $s^D$, then $s^D$ is the unique dominant strategy equilibrium.
\resolespaco

% ---------------------------------------------------------------------------
\exercicio[2]{IDSDS and Rationality}{Exam level}{Muñoz-Garcia \& Toro-Gonzalez (2019), \emph{Strategy
and Game Theory: Practice Exercises with Answers}, 2nd ed., Ch.~1 ``Dominance Solvable Games'',
Exercise 12 ``IDSDS and Rationality''$^{\mathrm{A}}$, p.~23 [PDF 36]. Worked answer in the book,
right after the statement. Reproduced verbatim (statement only).}

Consider the following $3 \times 3$ matrix (Fig.~1.30):
\jogo{Player 1}{Player 2}{\begin{tabular}{c|c|c|c|}
\hd{} & \hd{$F$} & \hd{$C$} & \hd{$B$}\\\cline{2-4}
$F$ & 0, 5 & 2, 4 & 2, 3\\\cline{2-4}
$C$ & 2, 3 & 0, 5 & 3, 2\\\cline{2-4}
$B$ & 5, 0 & 2, 4 & 2, 3\\\cline{2-4}
\end{tabular}}
\vspace{-6pt}\begin{center}\small\textit{Fig.~1.30} Normal-form representation of a three-players game\end{center}
\begin{partes}
\item Does any player have a strictly dominant strategy?
\item Does any player have a strictly dominated strategy?
\item Use your results from part (b) to delete players' strictly dominated strategies assuming
rationality, but not assuming common knowledge of rationality. [Hint: This means that every player
can delete the strategies he finds to be strictly dominated but cannot put himself in the shoes of
his opponent to delete the strategies his opponent finds to be strictly dominated.]
\item Use your results from part (b) to delete players' strictly dominated strategies assuming
common knowledge of rationality. [Hint: This means that players can put themselves in the shoes of
their opponents to delete strategies each other finds to be strictly dominated, as in the standard
approach of IDSDS.]
\end{partes}
\resolespaco

% ---------------------------------------------------------------------------
\exercicio[3]{Campaigning}{Exam level}{Tadelis (2013), \emph{Game Theory: An Introduction}, Ch.~4
``Rationality and Common Knowledge'', Exercise 4.7 ``Campaigning'', pp.~77--78 [PDF 96--97]. No
answer in the book. Reproduced verbatim.}

\rot{4.7}\textbf{Campaigning:} Two candidates, 1 and 2, are running for office. Each has one of
three choices in running his campaign: focus on the positive aspects of one's own platform (call
this a positive campaign [$P$]), focus on the positive aspects of one's own platform while
attacking one's opponent's campaign (call this a balanced campaign [$B$]), or finally focus only on
attacking one's opponent (call this a negative campaign [$N$]). All a candidate cares about is the
probability of winning, so assume that if a candidate expects to win with probability
$\pi \in [0, 1]$, then his payoff is $\pi$. The probability that a candidate wins depends on his
choice of campaign and his opponent's choice. The probabilities of winning are given as follows:

If both choose the same campaign then each wins with probability 0.5.

If candidate $i$ uses a positive campaign while $j \neq i$ uses a balanced one then $i$ loses for
sure.

If candidate $i$ uses a positive campaign while $j \neq i$ uses a negative one then $i$ wins with
probability 0.3.

If candidate $i$ uses a negative campaign while $j \neq i$ uses a balanced one then $i$ wins with
probability 0.6.
\begin{partes}
\item Model this story as a normal-form game. (It suffices to be specific about the payoff function
of one player and to explain how the other player's payoff function is different and why.)
\item Write down the game in matrix form.
\item What happens at each stage of elimination of strictly dominated strategies? Will this
procedure lead to a clear prediction?
\end{partes}
\resolespaco

% ---------------------------------------------------------------------------
\exercicio[3]{Discrete First-Price Auction}{Exam level}{Tadelis (2013), \emph{Game Theory: An
Introduction}, Ch.~4 ``Rationality and Common Knowledge'', Exercise 4.3 ``Discrete First-Price
Auction'', p.~76 [PDF 95]. No answer in the book. Reproduced verbatim.}

\rot{4.3}\textbf{Discrete First-Price Auction:} An item is up for auction. Player 1 values the item
at 3 while player 2 values the item at 5. Each player can bid either 0, 1, or 2. If player $i$ bids
more than player $j$ then $i$ wins the good and pays his bid, while the loser does not pay. If both
players bid the same amount then a coin is tossed to determine who the winner is, and the winner
gets the good and pays his bid while the loser pays nothing.
\begin{partes}
\item Write down the game in matrix form.
\item Does any player have a strictly dominated strategy?
\item Which strategies survive IESDS?
\end{partes}
\resolespaco

% ---------------------------------------------------------------------------
\exercicio[4]{Roommates}{Exam level}{Tadelis (2013), \emph{Game Theory: An Introduction}, Ch.~4
``Rationality and Common Knowledge'', Exercise 4.6 ``Roommates'', p.~77 [PDF 96]. No answer in the
book. Reproduced verbatim.}

\rot{4.6}\textbf{Roommates:} Two roommates each need to choose to clean their apartment, and each
can choose an amount of time $t_i \ge 0$ to clean. If their choices are $t_i$ and $t_j$, then
player $i$'s payoff is given by $(10 - t_j)t_i - t_i^2$. (This payoff function implies that the more
one roommate cleans, the less valuable is cleaning for the other roommate.)
\begin{partes}
\item What is the best response correspondence of each player $i$?
\item Which choices survive one round of IESDS?
\item Which choices survive IESDS?
\end{partes}
\resolespaco

% ---------------------------------------------------------------------------
\exercicio[4]{Order of deletion}{Advanced}{Mas-Colell, Whinston \& Green (1995), \emph{Microeconomic
Theory}, Ch.~8 ``Simultaneous-Move Games'', Exercises, Exercise 8.B.4$^{\mathrm{C}}$, p.~262 [PDF
142, left half]. Answer in the MWG solutions manual (\texttt{Livros/}). Reproduced verbatim.}

\rot{8.B.4$^{\mathrm{C}}$}Show that the order of deletion does not matter for the set of strategies
surviving a process of iterated deletion of strictly dominated strategies.
\resolespaco

% ---------------------------------------------------------------------------
\exercicio[5]{Cournot by successive elimination}{Advanced}{Mas-Colell, Whinston \& Green (1995),
\emph{Microeconomic Theory}, Ch.~8 ``Simultaneous-Move Games'', Exercises, Exercise
8.B.5$^{\mathrm{C}}$, p.~262 [PDF 142, left half]. Answer in the MWG solutions manual
(\texttt{Livros/}; 8.B.5 is there). Reproduced verbatim. The pointer to Chapter~12 is context; the
statement defines the whole model.}

\rot{8.B.5$^{\mathrm{C}}$}Consider the Cournot duopoly model (discussed extensively in Chapter 12)
in which two firms, 1 and 2, simultaneously choose the quantities they will sell on the market,
$q_1$ and $q_2$. The price each receives for each unit given these quantities is
$P(q_1, q_2) = a - b(q_1 + q_2)$. Their costs are $c$ per unit sold.

\textbf{(a)} Argue that successive elimination of strictly dominated strategies yields a unique
prediction in this game.

\textbf{(b)} Would this be true if there were three firms instead of two?
\resolespaco
